\input{../tex-inputs/latex-leseansicht-vorspann}

\section[Carl von Torresani an Arthur Schnitzler, 8. 5. 1893]{L04330 Carl von Torresani an Arthur Schnitzler, 8. 5. 1893}
\nopagebreak\mylabel{L04330v}
\rehead{ }\normalsize\beginnumbering\briefempfaengerindex{Schnitzler, Arthur@\textsc{Schnitzler, Arthur}!zzzTorresani-Lanzenfeld, Carl von@\emph{von Carl von Torresani-Lanzenfeld}!1893-05-081@{8. 5. 1893}|(be}
\toendnotes[C]{\smallbreak\pagebreak[2]}
\correspDesc{Versand  durch Carl von Torresani am 8. 5. 1893 in Venedig
\newline{}Erhalt  durch Arthur Schnitzler im Zeitraum [9. 5. 1893 – 13. 5. 1893?] in Wien}\toendnotes[C]{\smallbreak}
\Standort{TMW, Schn 4/86/1.}
\physDesc{Brief, 1 Blatt, 2 Seiten, 838 Zeichen
\newline{}Handschrift: schwarze Tinte, deutsche Kurrent}\toendnotes[C]{\smallbreak}
\pstart
           {\pb}Venedig\oindex{Venedig@\textbf{Venedig}|pw}{ }8/5 93\pend
           
\pstart{}Hochgeſchätzter Herr Doktor!\pend\vspace{0.5em}
\pstart
           Zu meinem lieben bedauern las ich in den Zeitungen die Nachricht von dem
               unerſetzlichen Verluſte\pwindex{Schnitzler, Johann 10.\,4.\,1835 Nagykanizsa – 2.\,5.\,1893 Wien@\textsc{Schnitzler, Johann} (10.\,4.\,1835 Nagykanizsa – 2.\,5.\,1893 Wien), \emph{Laryngologe}|pwv} der
               Sie betroffen hat; geſtatten Sie mir, Ihnen Sie in aufrichtiger und herzlicher Weise
               meines Mitgefühles zu verſichern. Es{ }ſind nun{ }ſchon einige Tage vergangen, Sie werden
               das Bitterſte bereits verwunden haben u.{ }ſich in philoſophiſcher Weiſe zu tröſten
                  wiſ{[}ſ{]}en. Aber in{ }ſolchen {\pb}Fällen muß man mit Romeo\pwindex{Shakespeare, William 23.\,4.\,1564? Stratford-upon-Avon – 3.\,5.\,1616 ebd.@\textsc{Shakespeare, William} (23.\,4.\,1564? Stratford-upon-Avon – 3.\,5.\,1616 ebd.), \emph{Schauspieler, Dramatiker}!Romeo und Julia. Trauerspiel in fünf Akten@\strich\emph{Romeo und Julia. Trauerspiel in fünf Akten}|pwv}{ }ſagen: »\label{K_L04330-1v}\edtext{Hängt die Philosophie.}{\lemma{\textnormal{\emph{Hängt die Philosophie.}}}\Cendnote{\textnormal{»Hängt die Philosophie! Kann sie nicht schaffen
                        eine Julia, Aufheben eines Fürsten Urteilsspruch, Verpflanzen eine Stadt: so
                        hilft sie nicht, So taugt sie nicht; so rede länger nicht!« William Shakespeare\pwindex{Shakespeare, William 23.\,4.\,1564? Stratford-upon-Avon – 3.\,5.\,1616 ebd.@\textsc{Shakespeare, William} (23.\,4.\,1564? Stratford-upon-Avon – 3.\,5.\,1616 ebd.), \emph{Schauspieler, Dramatiker}|pwk}: \emph{Romeo und Julia}\pwindex{Shakespeare, William 23.\,4.\,1564? Stratford-upon-Avon – 3.\,5.\,1616 ebd.@\textsc{Shakespeare, William} (23.\,4.\,1564? Stratford-upon-Avon – 3.\,5.\,1616 ebd.), \emph{Schauspieler, Dramatiker}!Romeo und Julia. Trauerspiel in fünf Akten@\strich\emph{Romeo und Julia. Trauerspiel in fünf Akten}|pwk}, übersetzt von August Wilhelm von Schlegel\pwindex{Schlegel, August Wilhelm von 5.\,9.\,1767 Hannover – 12.\,5.\,1845 Bonn@\textsc{Schlegel, August Wilhelm von} (5.\,9.\,1767 Hannover – 12.\,5.\,1845 Bonn), \emph{Schriftsteller, Kritiker, Übersetzer}|pwk}, 3. Akt,
                     3. Szene.}}}\label{K_L04330-1}\pwindex{Shakespeare, William 23.\,4.\,1564? Stratford-upon-Avon – 3.\,5.\,1616 ebd.@\textsc{Shakespeare, William} (23.\,4.\,1564? Stratford-upon-Avon – 3.\,5.\,1616 ebd.), \emph{Schauspieler, Dramatiker}!Romeo und Julia. Trauerspiel in fünf Akten@\strich\emph{Romeo und Julia. Trauerspiel in fünf Akten}|pwv}« Sie tröſtet immer nur die{ }ſie nicht brauchen, und
               ein wahrer Schmerz muß doch in{ }ſich{ }ſelbſt erſchöpfen. Möge die Dauer bei ihnen eine
               nichtige{ }ſein.\pend
           
\pstart
           In herzlicher Freundſchaft u. Hochſchätzung{\\[\baselineskip]} Ihr ergebenſter{\\[\baselineskip]}\spacefill\mbox{Torresani}\pend
           \leftskip=0em{}
\pstart
           \noindent{}Für alle Fälle theile ich Ihnen eine Adreſſe mit, mittels welcher Briefe mich zu
                  allen Zeiten treffen: »\label{K_L04330-2v}\edtext{H. u. G. Adv.}{\lemma{\textnormal{\emph{H. u. G. Adv.}}}\Cendnote{\textnormal{Hof- und Gerichtsadvokat}}}\label{K_L04330-2} D\textsuperscript{r}{ }Klemperer\pwindex{Klemperer, Alois 21.\,6.\,1846 Prag – 8.\,4.\,1910 Wien@\textsc{Klemperer, Alois} (21.\,6.\,1846 Prag – 8.\,4.\,1910 Wien), \emph{Rechtsanwalt}|pw}, I Tuchlauben 7 Wien\oindex{Wien@\textbf{Wien}!I., Innere Stadt@\textbf{I., Innere Stadt}!Tuchlauben 7@\textbf{Tuchlauben 7}, \emph{Wohngebäude}|pw}.\pend
           \selectlanguage{ngerman}\endnumbering\briefempfaengerindex{Schnitzler, Arthur@\textsc{Schnitzler, Arthur}!zzzTorresani-Lanzenfeld, Carl von@\emph{von Carl von Torresani-Lanzenfeld}!1893-05-081@{8. 5. 1893}|)be}\mylabel{L04330h}
\begin{anhang}
\end{anhang}\newcommand{\dateiname}{L04330}\newcommand{\titel}{Carl von Torresani an Arthur Schnitzler, 8. 5. 1893}\newcommand{\editorInnen}{Herausgegeben von Jahnke, SelmaMüller, Martin Anton}\input{../tex-inputs/latex-leseansicht-abspann}
